\chapter{STEP import and export}\label{ch:step}

STEP (ISO 10303, ``Standard for the Exchange of Product model data'') is the neutral file
format that practically every CAD system reads and writes. \NF{} contains its own STEP writer
and reader in the C++ core. Neither needs OpenCASCADE, gmsh or any other library. Both work
directly on the kernel's data structures: shells, trimmed faces and NURBS geometry. This
chapter explains how to use both directions, every option, what ends up in the file, how
foreign files are translated and where the limits are.

There are two ways to produce STEP files from \NF{}:

\begin{description}[leftmargin=1.5em, style=nextline]
  \item[Native writer (\apiref{write_step}, \apiref{to_step})] Part of the kernel. It writes the
    B-spline data unchanged and keeps the topology of the shell exactly as it is: shared edges,
    vertices and seams. This is the recommended route.
  \item[gmsh route (\apiref{mesh.export_step})] Hands the geometry to gmsh's OpenCASCADE kernel,
    which then writes STEP, BREP or IGES (\cref{sec:step-gmsh}). It needs the optional gmsh
    package and is mainly useful for BREP and IGES output, or for writing loose surfaces that
    do not form a shell.
\end{description}

Reading is always done by the native reader (\apiref{read_step}, \apiref{from_step} and their
\code{_named} variants).

\section{Writing STEP files}\label{sec:step-write}

\subsection{The basic call}

\apiref{write_step} takes a path and either a single \code{Shell} or a list of shells:

\begin{lstlisting}[style=python]
import nurbforge as nf

part = nf.boolean(nf.make_box((-2, -2, 0), (4, 4, 2)),
                  nf.make_cylinder((0, 0, -1), (0, 0, 1), 0.5, 4), nf.BooleanOp.Difference)
nf.write_step("block.step", part)                     # one shell
nf.write_step("parts.step", [part, other], unit="m")  # several shells
text = nf.to_step(part)                               # the file contents as str
\end{lstlisting}

\apiref{to_step} accepts exactly the same arguments, without the path, and returns the
complete file contents as a Python string. This is useful for storing models in a database,
sending them over a network or inspecting them in tests. \code{write_step(path, ...)} and
\code{open(path).read() == to_step(...)} give identical text if all options, including the
time stamp, are equal.

Only shells can be written. A single surface or trimmed surface is first wrapped into a shell:
\code{nf.Shell.sew([nf.TrimmedSurface(surface)])}. The result is then written as an open
surface model.

\subsection{Options}

All options are keyword arguments with defaults. They are identical for \code{write_step} and
\code{to_step}:

\begin{longtable}{@{}p{0.2\textwidth}p{0.15\textwidth}p{0.57\textwidth}@{}}
\toprule
\textbf{Option} & \textbf{Default} & \textbf{Meaning} \\
\midrule\endhead
\code{unit} & \code{"mm"} & Length unit declared in the file: \code{"mm"}, \code{"cm"},
  \code{"m"}, \code{"inch"} or \code{"um"} (micrometre, written as
  \code{SI_UNIT(.MICRO.,.METRE.)}). Coordinates are \emph{not} converted. The unit only states
  what one model unit means, and importing programs convert on their side. A cube with edge
  length 2 written with \code{unit="m"} therefore appears in OpenCASCADE, which works in
  millimetres, with a volume of $8\cdot10^{9}$. \\
\code{name} & \code{"nurbforge"} & Name of the \code{PRODUCT} and the name field of the
  \code{FILE_NAME} header entry. CAD systems show it as the part name. \\
\code{author} & \code{""} & Author field of \code{FILE_NAME}. \\
\code{organization} & \code{""} & Organisation field of \code{FILE_NAME}. \\
\code{timestamp} & current UTC time & ISO 8601 time stamp in \code{FILE_NAME}. Set it
  explicitly (for example \code{"2026-01-01T00:00:00"}) to get byte-identical, reproducible
  files, for instance under version control or in regression tests. \\
\code{uncertainty} & $10^{-6}$ & Value of \code{UNCERTAINTY_MEASURE_WITH_UNIT} in model units.
  Importers such as OpenCASCADE use it as their initial vertex and edge tolerance. It should be
  at least as large as the sewing tolerance of the shells. \\
\code{names} & \code{[]} & One name per shell. Each name becomes the name of that shell's
  \code{MANIFOLD_SOLID_BREP} or \code{BREP_WITH_VOIDS}, and of the \code{CLOSED_SHELL} or
  \code{OPEN_SHELL}. \apiref{read_step_named} returns these names, so applications can assign
  materials, roles or boundary conditions per body. If given, the list must have exactly one
  entry per shell. \\
\bottomrule
\end{longtable}

\begin{note}[Units in the C++ API]
In C++ the options are fields of \code{nf::StepOptions}, and the unit is the enumeration
\code{nf::StepLengthUnit} (\code{Millimetre}, \code{Centimetre}, \code{Metre}, \code{Inch},
\code{Micrometre}). The shell names are \code{StepOptions::shell_names}.
\end{note}

\subsection{A complete example}

The following script writes a drilled block with all header options set. It then compares the
file with the string variant and counts the main entities. Finally it writes three named
shells into one file, one of them open, and reads the names back.

\example{14_step_export.py}

Output:

\begin{lstlisting}[style=text]
part: 7 faces, closed=True, volume=30.429204
file and string identical: True
MANIFOLD_SOLID_BREP          1
CLOSED_SHELL                 1
ADVANCED_FACE                7
EDGE_CURVE                   17
VERTEX_POINT                 12
B_SPLINE_SURFACE_WITH_KNOTS  7
B_SPLINE_CURVE_WITH_KNOTS    17
schema: AUTOMOTIVE_DESIGN { 1 0 10303 214 1 1 1 1 }
open shell closed: False
  Block  faces= 7 closed=True
  Plug   faces= 3 closed=True
  Cover  faces= 5 closed=False
\end{lstlisting}

The block has seven faces: four sides, top and bottom with a hole each, and the bore.
These seven faces have 17 edges in total; every one becomes exactly one \code{EDGE_CURVE},
shared by the two faces that meet there. The schema is AP214 (\code{AUTOMOTIVE_DESIGN}), which
every common CAD system reads.

\subsection{What is written: the file structure}\label{sec:step-structure}

Knowing the structure helps when a file is rejected by another program, or when you post-process
files yourself.

\paragraph{Header.} \code{FILE_DESCRIPTION}, \code{FILE_NAME} (name, time stamp, author,
organisation) and \code{FILE_SCHEMA(('AUTOMOTIVE_DESIGN { 1 0 10303 214 1 1 1 1 }'))}.

\paragraph{Product structure.} \code{APPLICATION_CONTEXT},
\code{APPLICATION_PROTOCOL_DEFINITION}, \code{PRODUCT}, \code{PRODUCT_DEFINITION_FORMATION},
\code{PRODUCT_DEFINITION}, \code{PRODUCT_DEFINITION_SHAPE} and
\code{SHAPE_DEFINITION_REPRESENTATION}. All shells of one call belong to one product. There is
no assembly structure.

\paragraph{Representation context.} \code{GEOMETRIC_REPRESENTATION_CONTEXT(3)} with the length
unit (an \code{SI_UNIT} with prefix, or \code{CONVERSION_BASED_UNIT('INCH', ...)}), radian,
steradian and the uncertainty.

{\raggedright \paragraph{Closed shells (solids).} Each closed shell becomes a
\code{MANIFOLD_SOLID_BREP(CLOSED_SHELL(...))} inside an
\code{ADVANCED_BREP_SHAPE_REPRESENTATION}.\par}
\begin{itemize}
  \item A shell that is not consistently oriented is oriented first, with \apiref{Shell.oriented}.
  \item Every connected component is classified by its signed volume. A component whose
    normals point outwards is a solid of its own. A component whose normals point inwards is a
    void; an example is the inner sphere of a cube with a spherical cavity. A void becomes an
    \code{ORIENTED_CLOSED_SHELL(..., .F.)} of the smallest solid that contains it, and that
    solid is then written as a \code{BREP_WITH_VOIDS}.
\end{itemize}

{\raggedright \paragraph{Open shells (surface models).} Each open shell becomes a
\code{SHELL_BASED_SURFACE_MODEL(OPEN_SHELL(...))} inside a
\code{MANIFOLD_SURFACE_SHAPE_REPRESENTATION}. If one call contains both closed and open
shells, all items go into one plain \code{SHAPE_REPRESENTATION}. Empty shells are skipped.\par}

\paragraph{Faces.} Every face is an \code{ADVANCED_FACE(bounds, surface, same_sense)}. The
outer trim loop becomes the \code{FACE_OUTER_BOUND} and every hole a \code{FACE_BOUND}. Each
bound is an \code{EDGE_LOOP} of \code{ORIENTED_EDGE}s. The face orientation is carried by the
parametrisation of the surface, because \apiref{TrimmedSurface.flipped} reparametrises the
surface instead of setting a flag. Therefore \code{same_sense} is \code{.T.} except where a
face must point against its surface normal, as in voids and inside-out components. There
\code{same_sense} and the bound orientations are \code{.F.}.

{\raggedright \paragraph{Edges and vertices.} Every edge of the shell (\code{ShellEdge}) is written once, as
\code{EDGE_CURVE(VERTEX_POINT, VERTEX_POINT, curve, .T.)}. All its uses refer to it through
\code{ORIENTED_EDGE}s, with orientation \code{.F.} where the loop runs against the curve
direction. Some special cases:\par}
\begin{itemize}
  \item The seam of a closed surface, for example the line where a cylinder closes, appears
    twice in the same loop with opposite orientations.
  \item The poles of revolved surfaces, such as the tips of a sphere, are degenerate boundary
    pieces and are not written as edges.
  \item Sewing can split an edge at a T-junction. The resulting pieces are ordered along their
    loop by chaining their vertices.
\end{itemize}

\paragraph{Geometry.} Surfaces become \code{B_SPLINE_SURFACE_WITH_KNOTS}, and curves become
\code{B_SPLINE_CURVE_WITH_KNOTS}. Knots are written as distinct values with multiplicities
(\code{knot_spec .UNSPECIFIED.}, \code{self_intersect .F.}), and the \code{closed} flags are
computed from the control net. Rational geometry, which includes every exact circle, cylinder,
sphere, cone and torus, needs the STEP ``complex instance'' notation:

\begin{lstlisting}[style=text]
(BOUNDED_SURFACE() B_SPLINE_SURFACE(pu,pv,net,.UNSPECIFIED.,uc,vc,.F.)
 B_SPLINE_SURFACE_WITH_KNOTS(mu,mv,ku,kv,.UNSPECIFIED.) GEOMETRIC_REPRESENTATION_ITEM()
 RATIONAL_B_SPLINE_SURFACE(weights) REPRESENTATION_ITEM('') SURFACE())
\end{lstlisting}

The partial entity types are listed in alphabetical order, each with only its own explicit
attributes, as the standard requires. Control points and weights are written with the
$u$ index outermost, which is \NF{}'s storage order.

\paragraph{Numbers.} Real numbers use the shortest representation that reads back to the same
double (\code{std::to_chars}). They always contain a decimal point (\code{1.}, \code{0.25},
\code{1.E-05}), so no precision is lost in a write--read round trip. Long instances are
wrapped after commas.

\subsection{Limitations of the writer}

\begin{itemize}
  \item \textbf{Edge curves.} The 3D edge curves are the trim curves mapped into space. They are
    exact where the mapping allows it:
    \begin{itemize}
      \item On planar faces, the image of a trim circle is a rational circle.
      \item Along iso-parameter lines the iso-curve of the surface is used, for example for the
        end circles and the seam of a cylinder.
    \end{itemize}
    Everywhere else the edge is a cubic fit within $1/100$ of the sewing tolerance. Conics are
    written as B-splines, never as \code{CIRCLE} or \code{ELLIPSE} entities.
  \item \textbf{No p-curves.} \code{SURFACE_CURVE} and \code{PCURVE} entities are not written.
    Readers project the edges onto the faces themselves, which OpenCASCADE does without
    complaint.
  \item \textbf{One product.} There is no assembly structure, and no colours, layers or names
    per face. Face labels (\cref{ch:solids}) are not written.
  \item \textbf{Edges between shells.} Faces are written per shell. An edge shared by two
    different shells is written once per shell.
  \item \textbf{gmsh import.} gmsh's \code{importShapes} drops free faces from a file that also
    contains solids. This also happens with files written by OpenCASCADE itself, so in gmsh a
    mixed file shows only its solids.
\end{itemize}

\subsection{Validation of the writer}

The test suite writes shells natively and re-imports them with OpenCASCADE through gmsh. It
then compares OpenCASCADE's volume with \apiref{Shell.volume}. Every import yields exactly one
solid. The remaining differences on curved faces come from OpenCASCADE's approximate mass
integration; planar faces agree to rounding.

\begin{center}
\begin{tabular}{@{}lrrr@{}}
\toprule
Shape & \code{Shell.volume()} & OCC volume & rel.\ difference \\
\midrule
Cube $2\times2\times2$ & 8 & 8.000000000000 & $2\cdot10^{-16}$ \\
Unit sphere (rational, seam) & 4.188790205 & 4.187243221 & $3.7\cdot10^{-4}$ \\
Union of two unit spheres (distance 1) & 7.068583460 & 7.067161594 & $2.0\cdot10^{-4}$ \\
Cube minus cylinder $r=0.3$ & 7.434513326 & 7.435302412 & $1.1\cdot10^{-4}$ \\
Cube minus sphere $r=0.5$ (void) & 7.476401224 & 7.476577524 & $2.4\cdot10^{-5}$ \\
\bottomrule
\end{tabular}
\end{center}

The C++ tests also check the file syntax:
\begin{itemize}
  \item entity counts, and that every reference resolves;
  \item each edge is used exactly twice, with opposite orientations, and seams are correct;
  \item the order of the partial types in complex entities;
  \item the syntax of every real number.
\end{itemize}

\section{Reading STEP files}\label{sec:step-read}

\subsection{The reader functions}

The native reader imports the B-rep solids and surface models of STEP files. Files may come
from other CAD systems (AP203, AP214 or AP242; the schema name is not checked) or from
\code{write_step}. There are four entry points, which differ only in their input and output:

\begin{center}
\begin{tabularx}{\textwidth}{@{}lX@{}}
\toprule
Function & Input $\to$ result \\
\midrule
\apiref{read_step}\code{(path, *, unit, tolerance, skip_unsupported)} & file $\to$
  \code{list[Shell]} \\
\apiref{from_step}\code{(text, *, ...)} & file contents (\code{str}) $\to$ \code{list[Shell]} \\
\apiref{read_step_named}\code{(path, *, ...)} & file $\to$ \code{list[(name, Shell)]} \\
\apiref{from_step_named}\code{(text, *, ...)} & contents $\to$ \code{list[(name, Shell)]} \\
\bottomrule
\end{tabularx}
\end{center}

The options are keyword-only:

\begin{longtable}{@{}p{0.22\textwidth}p{0.12\textwidth}p{0.58\textwidth}@{}}
\toprule
\textbf{Option} & \textbf{Default} & \textbf{Meaning} \\
\midrule\endhead
\code{unit} & \code{"mm"} & Target unit. All coordinates are converted from the unit declared
  in the file into this unit. Possible values are \code{"mm"}, \code{"cm"}, \code{"m"},
  \code{"inch"} and \code{"um"}. Files without a unit declaration are read as millimetres. \\
\code{tolerance} & $10^{-6}$ & Sewing tolerance in the \emph{target} unit. Faces are sewn with
  \code{Shell.sew(faces, tolerance)}, and trim curves are rebuilt to $\text{tolerance}/10$.
  Files from systems that work with coarse tolerances (``sloppy'' files, where edges do not lie
  exactly on their faces) may need a larger value, such as $10^{-5}$ or $10^{-4}$ mm. \\
\code{skip_unsupported} & \code{False} & By default an unsupported entity raises
  \code{GeometryError}. With \code{True}, the affected face is left out and the message is
  issued as a Python \code{UserWarning}. The shell then stays open where the face is missing. \\
\bottomrule
\end{longtable}

The returned shells are ordinary \NF{} objects: all queries, Booleans, sections and meshing
apply to them directly.

\paragraph{Which shells are returned.}
\begin{itemize}
  \item Each \code{MANIFOLD_SOLID_BREP} becomes one shell.
  \item A \code{BREP_WITH_VOIDS} becomes \emph{one} shell that holds the outer boundary and the
    voids. Void faces point into the void, as \apiref{Shell.oriented} expects, so
    \code{volume()} is the volume of the material.
  \item Every shell of a \code{SHELL_BASED_SURFACE_MODEL} becomes one shell, usually open.
\end{itemize}

\paragraph{Names.} \code{read_step_named} and \code{from_step_named} return \code{(name, Shell)}
pairs. The name is the solid name, as written with \code{names=[...]}, or the \code{PRODUCT}
name if the solid has none. In C++ the pairs are \code{nf::NamedShell}.

\paragraph{The file's unit (C++).} \code{StepReadOptions::file_unit_metres} receives the length
unit declared in the file, in metres. One use is to detect a model that was drawn in
micrometres but declared as millimetres.

{\raggedright \paragraph{Warnings (C++).} With \code{skip_unsupported = true}, C++ callers can pass a
\code{std::vector<std::string>*} in \code{StepReadOptions::warnings} to collect the messages.
Python reports them as \code{UserWarning}s instead.\par}

\subsection{Example: units, names and errors}

\example{14_step_import.py}

Output:

\begin{lstlisting}[style=text]
hollow: components=2, volume=7.476401224
read as m : volume = 7.4764 (expected 7.4764), closed=True
read as mm: volume = 7.4764e+09 (expected 7.4764e+09), closed=True
read as cm: volume = 7.4764e+06 (expected 7.4764e+06), closed=True
Pillar  volume = 1570.7963 um^3
Cap     volume = 904.7787 um^3
GeometryError: read_step: unsupported entity OFFSET_SURFACE (#23) (face #54)
skipped: 1 warning(s), 5 faces, closed=False
\end{lstlisting}

The cube with a spherical cavity is one shell with two components. It is written as a
\code{BREP_WITH_VOIDS} and comes back as one closed shell with the same volume. The volume
scales with the cube of the unit factor, because the coordinates are converted. The error
message names the entity, its instance number (\texttt{\#23}) and the face it belongs to
(\texttt{\#54}), so the problem can be found in the file directly.

\subsection{Supported entities}

\begin{longtable}{@{}p{0.16\textwidth}p{0.78\textwidth}@{}}
\toprule
\textbf{Kind} & \textbf{Entities} \\
\midrule\endhead
Part 21 syntax & Header and \code{DATA} sections, simple and complex instances
  (\texttt{\#n=(A(...) B(...) ...)}), strings with doubled quotes, enumerations, \texttt{\$},
  \code{*}, typed parameters such as \code{LENGTH_MEASURE(1.)}, lists, comments, line breaks
  anywhere. \\
Topology & \code{MANIFOLD_SOLID_BREP}, \code{BREP_WITH_VOIDS}, \code{SHELL_BASED_SURFACE_MODEL},
  \code{CLOSED_SHELL}, \code{OPEN_SHELL}, \code{ORIENTED_CLOSED_SHELL},
  \code{ORIENTED_OPEN_SHELL}, \code{ADVANCED_FACE} and \code{FACE_SURFACE} (with
  \code{same_sense}), \code{FACE_OUTER_BOUND}, \code{FACE_BOUND} (with orientation),
  \code{EDGE_LOOP}, \code{VERTEX_LOOP} (ignored), \code{POLY_LOOP}, \code{ORIENTED_EDGE},
  \code{EDGE_CURVE} (with \code{same_sense}), \code{VERTEX_POINT}. \\
Surfaces & \code{B_SPLINE_SURFACE_WITH_KNOTS}, \code{BEZIER_SURFACE}, \code{UNIFORM_SURFACE},
  \code{QUASI_UNIFORM_SURFACE} (also as rational complex instances); \code{PLANE},
  \code{CYLINDRICAL_SURFACE}, \code{CONICAL_SURFACE}, \code{SPHERICAL_SURFACE},
  \code{TOROIDAL_SURFACE}. \\
Curves & \code{B_SPLINE_CURVE_WITH_KNOTS}, \code{BEZIER_CURVE}, \code{UNIFORM_CURVE},
  \code{QUASI_UNIFORM_CURVE} (also rational), \code{LINE}, \code{CIRCLE}, \code{ELLIPSE},
  \code{HYPERBOLA}, \code{PARABOLA}, \code{POLYLINE}. \code{SURFACE_CURVE}, \code{SEAM_CURVE}
  and \code{INTERSECTION_CURVE} contribute their 3D curve. For \code{TRIMMED_CURVE} the basis
  curve is used, and the vertices bound the edge. \\
Placement & \code{CARTESIAN_POINT}, \code{DIRECTION}, \code{VECTOR},
  \code{AXIS2_PLACEMENT_3D}. \\
Units & \code{SI_UNIT} with any prefix, and \code{CONVERSION_BASED_UNIT} (for example
  \code{'INCH'} via \code{LENGTH_MEASURE_WITH_UNIT}), for length and plane angle. Plane angles
  are needed for the semi-angles of cones. Units are taken from the
  \code{GLOBAL_UNIT_ASSIGNED_CONTEXT}. \\
\bottomrule
\end{longtable}

Other entities found while translating a face raise \code{GeometryError} naming the entity and
instance; examples are \code{OFFSET_SURFACE} and \code{SURFACE_OF_REVOLUTION}. With
\code{skip_unsupported}, that face is skipped as described above. A \code{FACETED_BREP} is
reported the same way. Product structure, presentation (colours, layers), PMI and p-curves are
ignored.

\subsection{How geometry is translated}

\paragraph{B-splines} are taken over unchanged. Files with unclamped knot vectors are clamped
first.

\paragraph{Analytic surfaces} become exact rational NURBS patches. In each case the normal
$\mathbf S_u\times\mathbf S_v$ equals the STEP surface normal.
\begin{itemize}
  \item A plane becomes a rectangle around the face boundary, enlarged by 5\,\%.
  \item Cylinders and cones become a revolved line. Cones stop at the apex.
  \item A sphere becomes a revolved half circle, and a torus a revolved circle.
\end{itemize}
A closed surface is represented by a patch covering only part of the full turn, and the
angular range is chosen from the face boundary:
\begin{itemize}
  \item The largest angular gap that no boundary curve crosses lies outside the face,
    provided the boundary turns away from it (judged by the inward boundary normal). The
    patch covers the rest.
  \item If there is no such gap, the patch covers the full turn. It starts at a seam edge of
    the face, which is an edge used twice in the same loop, so the seams of the file and of
    the patch coincide.
\end{itemize}

\paragraph{Edges} are built from their vertices.
\begin{itemize}
  \item Lines and conic arcs are exact. Circles and ellipses become rational arcs, hyperbolas
    and parabolas rational quadratic B\'ezier segments.
  \item B-spline edges become the sub-curve between the vertex parameters. A closed curve is
    wrapped around its start point if necessary.
  \item The end points are set exactly to the vertices.
\end{itemize}

\subsection{How trim loops are rebuilt}

STEP faces are bounded by 3D edges. \NF{}'s trimmed surfaces need loops in the $(u,v)$
parameter space of the surface (\cref{ch:trimming}). The reader maps every edge use into
parameter space, in one of three ways:

\begin{enumerate}
  \item \textbf{Planar patches with affine parametrisation.} This covers the writer's
    rectangles and the patches built for \code{PLANE}. The map is exact: the inverse affine map
    is applied to the edge curve.
  \item \textbf{Edges on an iso-parameter line.} Examples are the end circles and seams of
    cylinders and cones, and the meridians and parallels of spheres and tori. These are exact
    straight segments in parameter space. An edge qualifies if all samples share the fixed
    parameter to $10^{-9}$ of the domain, the running parameter is monotone, and the
    iso-curve between neighbouring samples lies on the edge.
  \item \textbf{All other edges} are projected sample by sample. Each sample uses a local
    Newton step from the previous one, with a global closest-point search as fallback. The
    samples are then interpolated with the edge parameter. The image of the interpolant is
    checked against the 3D edge between the samples and refined until it lies within
    $\text{tolerance}/10$, plus the edge's own distance from the surface.
\end{enumerate}

Closed surfaces need special care at the seam and the poles:
\begin{itemize}
  \item A sample that lies on the seam takes the side of its nearest regular sample.
  \item A curve that lies \emph{on} the seam appears twice in its loop and takes the side given
    by the orientation: the face lies to the left of the loop, seen against the face normal.
  \item Curves are split where they cross the seam or pass through a pole.
  \item A point at a pole takes the free parameter of its neighbours. Consecutive curves that
    end at different parameters of the same pole are connected along the pole.
  \item A loop that runs once around a closed surface, such as a band on a cylinder or torus
    without a seam edge, or a spherical cap, is closed along the seam or through the pole.
\end{itemize}
The loop with the largest area in parameter space becomes the outer loop. Faces with
\code{same_sense .F.}, and faces of an \code{ORIENTED_CLOSED_SHELL(.F.)}, are flipped with
\apiref{TrimmedSurface.flipped}.

\subsection{Validation of the reader}

The C++ tests read back the writer's output and check face counts, closedness, orientation,
area and volume to within $10^{-6}$. The models are a cube, a sphere, a cylinder with lids,
the union of two spheres, a cube with a spherical void, an open shell, several shells in one
file, and files in metres and inches. The tests also read hand-written files in the style of
other CAD systems: a cube made of \code{PLANE} and \code{LINE}, a cylinder made of
\code{CYLINDRICAL_SURFACE} and \code{CIRCLE}, a half cylinder, a sphere bounded only by a
\code{VERTEX_LOOP}, and a hemisphere. The Python tests read files written by gmsh's
OpenCASCADE kernel, if gmsh is installed, and compare the volumes, computed with
\code{Shell.volume(1e-7)}:

\begin{center}
\begin{tabular}{@{}lrrr@{}}
\toprule
OCC model & volume read & reference & rel.\ diff. \\
\midrule
Box $10\times20\times30$ & 6000 & 6000 & 0 \\
Cylinder $r=5$, $h=20$ & 1570.796327 & $500\pi$ & $10^{-16}$ \\
Cone frustum (seam, two circles) & 1455.604596 & $1390\pi/3$ & $3\cdot10^{-13}$ \\
Cone with apex & 261.799388 & $250\pi/3$ & $5\cdot10^{-12}$ \\
Sphere $r=4$ (\code{VERTEX_LOOP} only) & 268.0825731 & $256\pi/3$ & $2\cdot10^{-16}$ \\
Torus $R=20$, $r=5$ (seams both ways) & 9869.604401 & $1000\pi^2$ & $2\cdot10^{-13}$ \\
Box minus cylinder & 717.256661 & $1000-90\pi$ & $10^{-15}$ \\
Box, all edges filleted ($r=2$) & 907.7049927 & OCC 907.7049927 & $10^{-13}$ \\
Cone cut along its axis (\code{HYPERBOLA}) & 65.6912391 & OCC 65.6912385 & $9\cdot10^{-9}$ \\
Sphere cut through the poles & 56.5486678 & $18\pi$ & $10^{-12}$ \\
Box with spherical void & 966.489678 & $1000-32\pi/3$ & $10^{-15}$ \\
\bottomrule
\end{tabular}
\end{center}

\subsection{Limitations of the reader}

\begin{itemize}
  \item \textbf{Assemblies.} Placement transformations (\code{MAPPED_ITEM},
    \code{ITEM_DEFINED_TRANSFORMATION}, \code{CONTEXT_DEPENDENT_SHAPE_REPRESENTATION}) are not
    applied. Every solid is returned in the coordinates of its own representation. Only one
    length unit per file is used: the first \code{GLOBAL_UNIT_ASSIGNED_CONTEXT}.
  \item \textbf{Unsupported geometry.} \code{SURFACE_OF_REVOLUTION},
    \code{SURFACE_OF_LINEAR_EXTRUSION}, \code{OFFSET_SURFACE},
    \code{RECTANGULAR_TRIMMED_SURFACE} and \code{CURVE_BOUNDED_SURFACE} are not supported.
    Neither are faceted B-reps, tori with minor radius $\ge$ major radius, wireframe and mesh
    geometry.
  \item \textbf{P-curves are not used.} Trim loops are always rebuilt from the 3D edges, so the
    edges must lie on their faces within the tolerance. Sloppy files may need a larger
    \code{tolerance}.
  \item \textbf{Seam crossings.} On a closed B-spline surface, a face whose boundary crosses the
    seam twice without a seam edge (a patch straddling the seam) is rejected. The surface is
    not reparametrised.
\end{itemize}

\begin{tip}[Debugging a file that does not import]
\begin{enumerate}
  \item Read the error message. It names the entity type, the instance number and the face.
    Open the file in a text editor and search for the instance, for example \texttt{\#23=}.
  \item Try \code{skip_unsupported=True} to see how many faces are affected. If only a few
    faces are missing, the open edges show where (\apiref{Shell.free_edges}).
  \item If edges do not meet, increase \code{tolerance}, for example to \code{1e-5} or
    \code{1e-4} in the target unit.
  \item If the file has an assembly structure, export the part alone from the original CAD
    system, or apply the placements yourself with \apiref{Shell.transformed}.
\end{enumerate}
\end{tip}

\section{STEP, BREP and IGES through gmsh}\label{sec:step-gmsh}

\apiref{mesh.export_step} writes surfaces, trimmed surfaces and shells through gmsh's
OpenCASCADE kernel:

\begin{lstlisting}[style=python]
from nurbforge.mesh import export_step

export_step([part], "part.step")                 # format from the extension:
export_step([part], "part.brep")                 #   .step, .brep, .iges
export_step([lens_front, lens_back], "lens.iges") # loose surfaces are allowed too
export_step([part], "faces.step", solid=False)   # closed shells as faces only, no solid
\end{lstlisting}

Internally, each shell is handed to OpenCASCADE with the native writer and OCC's STEP reader,
as described in \cref{ch:meshing}. Surfaces are transferred exactly. OCC then writes the file
in its own way. \Cref{tab:step-routes} compares the two routes.

\begin{table}[htbp]
\centering
\begin{tabularx}{\textwidth}{@{}lXX@{}}
\toprule
& \textbf{Native} (\code{write_step}) & \textbf{gmsh / OCC} (\code{export_step}) \\
\midrule
Dependencies & none & gmsh Python package \\
Formats & STEP AP214 & STEP, BREP, IGES \\
Input & shells & shells, surfaces, trimmed surfaces \\
Topology & exactly the shell's edges and vertices & rebuilt by OCC \\
Header, units, names & all options of \cref{sec:step-write} & OCC defaults (mm) \\
Reproducible files & yes (fixed \code{timestamp}) & no \\
\bottomrule
\end{tabularx}
\caption{The two routes to a STEP file.}\label{tab:step-routes}
\end{table}

Use the native writer whenever a shell is to be written as STEP. Use the gmsh route for BREP
and IGES, or when you have loose surfaces that do not form a shell.

\section{STEP in the LLM interface}

The tool layer of \cref{ch:llm} has two STEP tools:
\begin{itemize}
  \item \toolref{export_step} writes objects with the native writer. Surfaces are sewn into a
    shell each, and the unit can be chosen.
  \item \toolref{import_step} reads a file and stores the shells as \code{name},
    \code{name_1}, \code{name_2}, and so on.
\end{itemize}
A model can therefore load an existing part, modify it with Booleans and write it out again,
all through whitelisted calls.
